% = PARAGRAPH 1 - Existing technical solutions for deepfake detection
% = (grounds the solution discussion in concrete tools)

There are various technical approaches to detect deepfakes. Forensic analysis involves examining the digital artifacts left by deepfake generation processes, such as inconsistencies in lighting, shadows, or facial movements. Authentication watermarking involves embedding a unique identifier into original content that can be verified later to confirm authenticity. However, these methods have limitations. As \citet{muam2026role} points out, deep learning and NLP techniques used for detection often lag behind the advancements in deepfake generation, making it a constant cat-and-mouse game between creators and detectors.


% = Paragraph 2 — Ethical frameworks for responsible AI development
% = (satisfies: "ethics / ethical aspects and frameworks")

Ethical frameworks for responsible AI development, such as those outlined by \citet{laws14060083}, emphasize principles like accountability, transparency, and fairness. In the context of deepfakes used in conflict, these principles raise critical questions about who is responsible for the harm caused. If a deepfake video incites violence or spreads misinformation, should the creators of the deepfake be held accountable? What about platforms that host such content, or do not take sufficient measures to fight misinformation? Sole responsibility is difficult to assign in the digital age, where content can be easily shared and manipulated.


% = Paragraph 3 — Regulatory responses
% = (satisfies: "regulation")

Certain regulatory responses have been proposed to address the challenges posed by deepfakes. The EU AI Act~\citep{EU_AI_Act_2024}, for example, aims to establish a legal framework for AI applications, including provisions for high-risk AI systems like deepfake generators. However, as \citet{laws14060083} argues, existing legal frameworks may not be sufficient to address the unique challenges of deepfakes in conflict contexts. International Humanitarian Law (IHL) may not adequately cover the use of deepfakes as a weapon of war, and there is a need for new regulations that specifically address the ethical and legal implications of AI-generated disinformation in conflict zones.

% = Paragraph 4 — Societal mechanisms
% = (satisfies: "society")

Beyond technical and regulatory solutions, media literacy programs can be effective to educate the general public about the existence and risks of deepfakes. NGOs and journalism standards organizations, such as the Content Authenticity Initiative, work to establish best practices for verifying content and maintaining journalistic integrity. International standards bodies like IEEE can also play a role in developing guidelines for responsible AI use. A collaboration between the technical, regulatory, and societal mechanisms is essential dampen the amplifying effect of societal factors like political polarization and weakened media literacy that make populations more vulnerable to deepfake disinformation.

% = Paragraph 5 — Personal reflection
% = (satisfies: "personal reflection")

It's my opinion that deepfakes are in fact an issue in that not everyone has the technical expertise to identify them. While I myself, and those I spend time with are mostly aware of deepfakes and their risks, I worry about the general public's ability this new landscape of disinformation. With how polarised the world is, combined with the fact that every social media had a tailored algorithm that feeds us content we are likely to blindly agree with, I think deepfakes are a real threat to the political climate. I don't think any single solution is sufficient to address this issue. In my opinion, and that of many others, a multi-layered approach that combines technical detection methods, regulatory frameworks, and societal mechanisms is necessary.

